\chapter{Setpoint coordination}
\label{ch:coord}
The toy examples in Chapter~\ref{ch:toy} asked the supervisor to detect
faults. The research questions instead ask how well a generated supervisor
\emph{coordinates setpoints} under disturbances compared with a model-based
predictive controller, inside and beyond the development range. This chapter
rehearses that comparison on the quadruple-tank process before the industrial
study. All scores in the chapter are evaluated on the same three batteries,
so every number is comparable; lower is better.

\section{Test bed}
\label{sec:coord-testbed}
\paragraph{Process and objectives.} The plant is the quadruple-tank process
of Section~\ref{sec:four} with the same cross-paired PI loops
($K_p = 40$, $K_i = 0.3$), which hold the levels $h_1$ and $h_2$ of the
lower tanks at their setpoints. Table~\ref{tab:coord-analog} maps the
grinding-circuit objectives onto the process. The product is the production
rate $Q$, the summed outflow of the two lower tanks, which is
16.4~L/s at the nominal setpoints $h_1 = 0.30$~m and $h_2 = 0.35$~m. The
level $h_2$ must stay inside the band 0.25--0.45~m, and the upper levels
$h_3$ and $h_4$ must stay below 0.75~m. The supervisor moves only the two
setpoints, within 0.05--0.48~m. Two setpoints and one production target
leave one degree of freedom, which the supervisor must use to keep the
constraints satisfied as disturbances move the operating point.

\begin{table}[tb]
  \centering
  \caption{The grinding-circuit objectives and their analogs on the
    quadruple-tank process.}
  \label{tab:coord-analog}
  \begin{tabular}{@{}ll@{}}
    \toprule
    Grinding circuit & Quadruple-tank analog \\
    \midrule
    P80 on target                  & production rate $Q$ on target \\
    Sump level stable              & $h_2$ inside a band \\
    Circulating load within limits & $h_3$ and $h_4$ below a limit \\
    Fresh feed flow variation      & draw-off at tank 1 or 2, extra inflow at tank 2 \\
    Ore density variation          & loss of pump gain $k_1$ or $k_2$ \\
    Percentage solids variation    & shift of valve split $\gamma_1$ or $\gamma_2$ \\
    Setpoint smoothness            & total variation of the setpoints \\
    \bottomrule
  \end{tabular}
\end{table}

\paragraph{Scenarios.} Each 1\,200~s episode starts in steady operation, as
a running plant would, and is integrated with the explicit Euler method at a
1~s step. A disturbance of one of three kinds (feed, pump or split, as in
Table~\ref{tab:coord-analog}) or a combination of two of them at three
quarters of the amplitude starts between 150 and 300~s, as a step, a ramp
over 200--400~s or a sine with a period of 200--500~s. The development
amplitudes were calibrated on the steady-state model so that some setpoint
pair still meets the production target with every constraint satisfied,
while the upper part of each range breaks the fixed operating recipe; the
amplitudes beyond the development range start past that edge of
solvability, so some constraint violation there cannot be avoided. Two
further nominal scenarios step the production target up or down by 5--10\%
and back. With three scenarios per cell, the development battery has 42
scenarios (12 cells of disturbance kind and pattern, plus the nominal
scenarios), the held-out battery has 42 scenarios of the same types with
other seeds, and the battery beyond the development range has 36. The
beyond-range battery is never shown to the LLM.

\paragraph{Supervisor interface.} The supervisor is called every 10~s as
\begin{center}
\texttt{supervise(telemetry\_window, active\_setpoints, objectives)}
\end{center}
and returns a diagnosis and the two new setpoints; there are no anomaly
flags. The window holds the last 600 one-second samples of the four levels,
both pump voltages, the production rate, and the setpoints and production
target that were active at each sample, as a distributed control system logs
them. The objectives are the production target, the band, the upper limit and
the setpoint limits. The supervisor is stateless, as before, so the window is
its only memory. An earlier 50~s window was too short for this: the
supervisor could not tell a sine disturbance from a ramp or a step, and it
could not see its own setpoint changes or a target change settling (about
70~s to within 2\%). The first generated supervisors read that transient as
a disturbance and chattered their setpoints.

\paragraph{Score.} For a scenario $s$,
\begin{equation}
  J_s = \mathrm{IAE}_{Q,s} + 2\,(B_s + U_s) + 10\,V_s + 100\,\mathrm{TV}_s
        + 10^3\,E_s ,
  \label{eq:coord-score}
\end{equation}
where $\mathrm{IAE}_{Q,s}$ is the integrated absolute production error in
litres, $B_s$ and $U_s$ the number of seconds with $h_2$ outside its band and
with an upper level above its limit, $V_s$ the number of seconds outside the
safety bounds of the lower levels, $\mathrm{TV}_s$ the total variation of the
two setpoints in metres and $E_s$ the number of exceptions and timeouts. The
weights are provisional: one litre off target costs as much as half a second
of constraint violation or a centimetre of setpoint travel. The recovery
time, the time until production settles within 2\% of the target after an
event, is reported but not scored. No controller in this chapter violated
the safety bounds or raised an exception.

\section{Baselines}
\label{sec:coord-baselines}
\paragraph{Fixed recipe.} The seed supervisor keeps a fixed operating
recipe: when the production target changes, it scales both setpoints with
the nominal square-root relation $Q \propto \sqrt{h}$, and it ignores
disturbances. This is what the PI loops achieve on their own.

\paragraph{MPC.} The model-based predictive baseline simulates candidate
setpoint pairs over a 300~s horizon with an internal model of the plant and
the PI loops, and applies the pair with the lowest predicted cost: squared
relative production error, squared band and upper-limit violations, pump
voltages outside 1.5--11.5~V, and setpoint moves. The candidates are a
coarse $9 \times 9$ grid of moves around the current setpoints and then a
finer $5 \times 5$ grid around the best one. The disturbance is estimated
online as an additive inflow per tank, corrected after each decision by half
the difference between the measured and the predicted state, which is the
bias update common in industrial MPC \cite{qin2003}. The \emph{oracle MPC}
instead uses the true plant parameters and disturbance at each decision,
held constant over the horizon, and bounds what a predictive controller with
this horizon and cost can reach. Unlike the generated supervisors, both are
stateful. They stand in for the model-based predictive benchmark of the
research questions; they are not the TD-MPC controller used at Weir.

\section{Changes to the training loop}
\label{sec:coord-loop}
The loop of Chapter~\ref{ch:method} was reused with the score
\eqref{eq:coord-score}, the same security sandbox, deepseek-flash and $K = 3$
candidates per generation. Four changes were made, each after a run that stalled:
\begin{itemize}
\item \emph{Regression guard by disturbance type, with a schedule.} The
  per-scenario guard of Section~\ref{sec:scoring} rejected almost every
  candidate; in the first 600~s run every candidate that beat the seed (the
  best scored 258 against 533) was rejected, so the prompt kept showing only
  the seed. Now the mean score of each disturbance type may get worse by at
  most the larger of 10 points and a relative tolerance, and no single
  scenario by more than the larger of an absolute and a relative cap. The
  tolerances are loose in generations 1--2 (25\%, cap 300 points or
  150\%), tighter in generations 3--4 (15\%, cap 200 points or 100\%) and
  tightest from generation 5 (10\%, cap 100 points or 50\%), and the battery
  average must always improve. A candidate that is much better overall but
  worse on one scenario can thus be promoted early, and its weak scenario
  then appears in the next prompt's decision traces.
\item \emph{Failure traces of rejected candidates.} The prompt also shows,
  for up to two non-promoted candidates of the previous generation, the
  decision trace of the scenario each lost most on.
\item \emph{One targeted change per candidate.} Without this rule,
  candidates rewrote the supervisor, keeping 16--65\% of the champion's lines,
  and each fixed one scenario while breaking others. With it they kept
  83--99\%.
\item \emph{A measured record instead of unverified lessons.} The lessons of
  Chapter~\ref{ch:method} are written by the LLM before its candidate is
  scored. In one run, all nine candidates made the same change (averaging the
  load estimate), which damped the response to sine disturbances and broke
  the response to steps, and the lessons saved from them recommended that
  change again. Lessons are now kept only from promoted candidates. Every
  other candidate is recorded by the trainer as one measured line (its
  change, its average against the champion, and the scenarios it gained and
  lost most on), and the prompt lists all of them for the run. A run seeded
  with another run's champion inherits that record.
\end{itemize}

\section{Runs}
\label{sec:coord-runs}
Table~\ref{tab:coord-runs} lists every run. Each row changed one part of the
loop from the row above. The first runs used the 50~s window and one
scenario per cell; their champions are re-scored on the current batteries.

\begin{table}[tb]
  \centering
  \caption{Coordination runs, in order. Gen.\ $\times$ $K$ is the number of
    generations and candidates per generation; cost is the estimated
    DeepSeek cost in US dollars. Scores are the champion's on the current
    batteries. All runs used low reasoning effort except the last.}
  \label{tab:coord-runs}
  \small\setlength{\tabcolsep}{4pt}
  \begin{tabular}{@{}>{\raggedright\arraybackslash}p{0.33\textwidth}ccrrrr@{}}
    \toprule
    Run (change from the row above) & Gen.\ $\times$ $K$ & Prom. & Cost
      & Dev. & Held-out & Beyond \\
    \midrule
    50~s window, per-scenario guard            & $2 \times 4$ & 2 & 0.31 & 252.1 & 241.3 & 1\,261 \\
    Same setup, three more runs                & $2 \times 4$ & 0 & 0.83 & 533.4 & 478.7 & 1\,583 \\
    Looser guard, run 1                        & $4 \times 4$ & 1 & 0.53 & 255.6 & 243.2 & 1\,190 \\
    \phantom{Looser guard,} run 2              & $4 \times 4$ & 2 & 0.35 & 226.3 & 242.2 & 1\,277 \\
    \phantom{Looser guard,} run 3              & $4 \times 4$ & 1 & 0.40 & 404.7 & 358.2 & 1\,390 \\
    600~s window, 3 scenarios per cell, guard by type
                                               & $3 \times 3$ & 0 & 0.32 & 533.4 & 478.7 & 1\,583 \\
    Guard schedule, failure traces             & $4 \times 3$ & 1 & 0.42 & 231.5 & 247.5 & 1\,339 \\
    One targeted change per candidate          & $3 \times 3$ & 0 & 0.11 & 231.5 & 247.5 & 1\,339 \\
    Measured record, lessons only when promoted & $3 \times 3$ & 0 & 0.21 & 231.5 & 247.5 & 1\,339 \\
    High reasoning effort, inherited record    & $4 \times 3$ & 4 & 0.63 & 206.3 & 214.9 & 1\,199 \\
    \bottomrule
  \end{tabular}
\end{table}

The guard schedule gave the first promotion with the 600~s window, 231.5
against the seed's 533.4, in its first generation. The next nine generations
at low reasoning effort, the rest of that run and two continuations from its
champion, promoted nothing, although the
measured record moved the candidates away from repeating failed changes and
on to estimators that accounted for oscillating disturbances (the best
scored 241.5). The same loop at high reasoning effort, seeded with that
champion and its record of nine measured attempts, promoted a candidate in
each of four generations: 230.3, 227.4, 212.4 and 206.3. The four changes
were, in order, a pump-gain estimate corrected for the flow stored in the
upper tanks, a feedback trim of the $h_1$ setpoint from the measured
production error, a predictive version of that trim, and a shorter hold after
the supervisor's own setpoint changes. A request at high effort produced
about 80\,000 output tokens, against 16\,000--40\,000 at low effort, and the
run cost about 0.63 US dollars. High effort and the inherited record were
introduced in the same run, so the run does not separate their effects.
All coordination runs together cost about 4.1 US dollars.

Two setups were run several times with identical settings. Of the four
runs with the 50~s window and the per-scenario guard, one promoted two
candidates and three promoted none, and the champions of the three runs with
the looser guard scored 255.6, 226.3 and 404.7 on the development battery. This is
a first indication, for contribution~C3, that independent runs of the same
setup can end far apart.

\section{Comparison with MPC}
\label{sec:coord-results}
Table~\ref{tab:coord-main} compares the best generated supervisor with the
baselines. On the development battery it is within 1.3\% of the MPC with an
estimated disturbance (206.3 against 203.7) and reduces the fixed recipe's
score by 61\%. On the held-out battery it is 11\% worse than the MPC (214.9
against 193.8). Its gap between the development and the held-out battery
grew from about 2 to about 9 points over the last two generations, while the
MPC, which was not tuned on the batteries, scored lower on the held-out
battery than on the development battery. This suggests that the last
promotions partly fit the development scenarios.

\begin{table}[tb]
  \centering
  \caption{Average score on the coordination batteries (lower is better).
    The beyond-range battery is never shown to the LLM.}
  \label{tab:coord-main}
  \begin{tabular}{@{}lrrr@{}}
    \toprule
    Controller & Development & Held-out & Beyond range \\
    \midrule
    Fixed recipe (seed)                 & 533.4 & 478.7 & 1\,583.0 \\
    Generated, low reasoning effort     & 231.5 & 247.5 & 1\,339.3 \\
    Generated, high reasoning effort    & 206.3 & 214.9 & 1\,199.2 \\
    MPC, estimated disturbance          & 203.7 & 193.8 & 1\,123.6 \\
    MPC, oracle                         & 192.9 & 191.1 & 1\,110.9 \\
    \bottomrule
  \end{tabular}
\end{table}

Table~\ref{tab:coord-metrics} splits the score into its terms, and
Table~\ref{tab:coord-kind} by disturbance type. On the development battery
the generated supervisor and the MPC are close on every term: production
error 155.7 against 153.1~L and constraint violation 17.8 against 19.3~s per
scenario, with somewhat more setpoint travel for the generated supervisor.
By disturbance type it is better than the MPC on feed disturbances (345.0
against 386.6), about equal on split disturbances, and worse on pump
disturbances, combined disturbances and production-target changes. For both
controllers, sine disturbances cost by far the most: 445.4 and 425.1 on
average against below 100 for steps and ramps.

\begin{table}[tb]
  \centering
  \caption{Score terms, averaged per scenario: production IAE, seconds of
    band or upper-limit violation, setpoint total variation and recovery
    time.}
  \label{tab:coord-metrics}
  \begin{tabular}{@{}llrrrr@{}}
    \toprule
    Battery & Controller & IAE (L) & Violation (s) & TV (m) & Recovery (s) \\
    \midrule
    Development & Fixed recipe & 199.1 & 165.8 & 0.03 & 288 \\
                & Generated    & 155.7 &  17.8 & 0.15 & 206 \\
                & MPC          & 153.1 &  19.3 & 0.12 & 193 \\
    \midrule
    Held-out    & Fixed recipe & 217.5 & 129.2 & 0.03 & 356 \\
                & Generated    & 181.0 &   9.1 & 0.16 & 260 \\
                & MPC          & 170.6 &   3.9 & 0.15 & 228 \\
    \midrule
    Beyond      & Fixed recipe & 767.3 & 407.9 & 0.00 & 675 \\
                & Generated    & 368.3 & 394.1 & 0.43 & 422 \\
                & MPC          & 454.1 & 312.7 & 0.44 & 577 \\
    \bottomrule
  \end{tabular}
\end{table}

\begin{table}[tb]
  \centering
  \caption{Average score by disturbance type on the development and the
    beyond-range battery. The beyond-range battery has no nominal
    scenarios.}
  \label{tab:coord-kind}
  \begin{tabular}{@{}lrrrrrr@{}}
    \toprule
    & \multicolumn{3}{c}{Development} & \multicolumn{3}{c}{Beyond range} \\
    \cmidrule(lr){2-4}\cmidrule(l){5-7}
    Type & Fixed & Generated & MPC & Fixed & Generated & MPC \\
    \midrule
    Nominal  & 251.8 & 230.1 & 213.0 &        --- &       --- &       --- \\
    Feed     & 950.1 & 345.0 & 386.6 & 2\,166.0 & 1\,525.1 & 1\,394.3 \\
    Pump     & 229.9 & 119.4 & 102.8 &    990.0 & 1\,213.3 &    931.2 \\
    Split    & 415.7 & 124.1 & 120.6 & 1\,539.0 &    769.0 &    871.8 \\
    Combined & 725.5 & 220.7 & 198.8 & 1\,637.0 & 1\,289.3 & 1\,297.0 \\
    \bottomrule
  \end{tabular}
\end{table}

\paragraph{Beyond the development range.} All controllers degrade sharply
beyond the development range, as intended by the calibration. The generated
supervisor's score rose by a factor of 5.8 (206.3 to 1\,199.2) and the MPC's
by 5.5 (203.7 to 1\,123.6), so the generated supervisor ends 6.7\% above the
MPC. The two degrade differently. The generated supervisor keeps production
closer to the target than the MPC (368.3 against 454.1~L) but violates the
constraints for longer (394.1 against 312.7~s), so it trades constraints for
production where the MPC does the opposite. By type, it does better than the
MPC on split disturbances and as well on combined ones, but on large pump-gain
losses it scores worse than the fixed recipe (1\,213.3 against 990.0). Why
has not yet been analyzed.
\TODO{Trace the beyond-range pump scenarios, and add a figure with one development and one beyond-range trajectory of both controllers.}

\section{The generated supervisor}
\label{sec:coord-champion}
The best supervisor is a single Python function of 261 lines. It works as a
steady-state optimizer with a feedback trim, a structure that the LLM arrived
at by itself from the seed's fixed recipe:
\begin{enumerate}
\item It estimates the extra inflow to each lower tank from a mass balance:
  the measured outflow minus the inflow that the upper tanks' outflows imply
  through the nominal valve splits.
\item It estimates each pump's gain from the flow into its upper tank,
  including the measured rate of change of the upper level, so that flow
  stored in an upper tank during an oscillation is not mistaken for a gain
  loss. The estimate is limited to 50--100\% of the nominal gain.
\item It divides the production target between the two lower tanks by
  trying 61 splits. For each it computes, from the steady-state model with
  the estimated disturbance, the required pump flows, the upper levels and
  the pump voltages, and it penalizes upper levels within 0.05~m of the limit,
  $h_2$ outside its band, saturated pumps and the size of the setpoint move.
  The resulting setpoints move at most 0.06~m per decision.
\item It trims the $h_1$ setpoint with the production error predicted 80~s
  ahead from the error's slope over the last 90~s. The prediction may
  reinforce the measured error but never change its sign, so steps and target
  changes can only be corrected faster, and the trim is limited by the
  remaining headroom of tank~3.
\item It holds its setpoints for 25~s after its own last setpoint change,
  unless the production target changed since, so that it does not react to
  its own transient.
\end{enumerate}
The code added by the promoted candidates is commented with the reasoning
behind it, in one place including the measured effect of an earlier
variant. The lessons kept from the four promoted
candidates are statements about the process, for example that ``a setpoint
change needs 1.5--2 minutes to reach the product, so a proportional error
trim corrects each swing after its trough''.
